\documentclass[11pt,a4paper]{article}

% shared/macros.tex
% Shared notation for all surface-partition math documents.
% Include with: % shared/macros.tex
% Shared notation for all surface-partition math documents.
% Include with: % shared/macros.tex
% Shared notation for all surface-partition math documents.
% Include with: \input{../shared/macros}

% -------------------------------------------------------------------------
% Packages (include here so each document only needs \input{../shared/macros})
% -------------------------------------------------------------------------
\usepackage{amsmath,amssymb,amsthm}
\usepackage{bm}           % \bm for bold math
\usepackage{mathtools}    % \coloneqq, dcases, etc.
\usepackage{booktabs}     % \toprule etc. in tables
\usepackage{hyperref}
\usepackage{geometry}
\usepackage{microtype}
\usepackage{parskip}
\usepackage{xcolor}
\usepackage{tcolorbox}
\usepackage{enumitem}
\usepackage{float}             % [H] float placement

\geometry{a4paper, margin=2.5cm}

% -------------------------------------------------------------------------
% Theorem-like environments
% -------------------------------------------------------------------------
\theoremstyle{definition}
\newtheorem{defn}{Definition}[section]
\newtheorem{remark}{Remark}[section]

\theoremstyle{plain}
\newtheorem{prop}{Proposition}[section]

% -------------------------------------------------------------------------
% Mesh and geometry
% -------------------------------------------------------------------------
\newcommand{\V}{\mathbf{V}}            % vertex array V[i] in R^dim
\newcommand{\F}{\mathbf{F}}            % face (triangle) array
\newcommand{\vone}[1]{v_{1,#1}}       % first vertex index of VP i
\newcommand{\vtwo}[1]{v_{2,#1}}       % second vertex index of VP i

% -------------------------------------------------------------------------
% Variable points
% -------------------------------------------------------------------------
\newcommand{\lam}{\lambda}             % scalar VP parameter
\newcommand{\Lam}{\boldsymbol{\lambda}}% full parameter vector
\newcommand{\vp}[1]{\mathbf{p}_{#1}}  % 3D position of VP i
\newcommand{\ed}[1]{\mathbf{d}_{#1}}  % edge direction d_i = V[v1_i] - V[v2_i]

% -------------------------------------------------------------------------
% Steiner / triple-point
% -------------------------------------------------------------------------
\newcommand{\Sv}{\mathbf{S}}           % Steiner (Fermat-Torricelli) point
\newcommand{\Ki}{K_{i}}               % projection matrix (I-n_i n_i^T)/r_i
\newcommand{\Mmat}{M}                 % sum of K_i matrices at a triple point

% -------------------------------------------------------------------------
% Normals and unit vectors
% -------------------------------------------------------------------------
\newcommand{\nhat}{\hat{\mathbf{n}}}  % unit normal to a triangle
\newcommand{\Dhat}{\hat{\boldsymbol{\Delta}}}  % unit segment direction

% -------------------------------------------------------------------------
% Operations
% -------------------------------------------------------------------------
\newcommand{\norm}[1]{\left\|#1\right\|}
\newcommand{\abs}[1]{\left|#1\right|}
\newcommand{\ip}[2]{\langle #1,\, #2 \rangle}

% Projection perpendicular to a unit vector \uhat
\newcommand{\Proj}[1]{\bigl(\mathbf{I} - #1 #1^\top\bigr)}

% argmax operator (winner-take-all assignment in Phase 1 documents)
\DeclareMathOperator*{\argmax}{arg\,max}
\DeclareMathOperator*{\argmin}{arg\,min}

% -------------------------------------------------------------------------
% Calculus shorthands
% -------------------------------------------------------------------------
\newcommand{\pd}[2]{\dfrac{\partial #1}{\partial #2}}
\newcommand{\pdd}[3]{\dfrac{\partial^2 #1}{\partial #2\,\partial #3}}
\newcommand{\grad}{\nabla}
\newcommand{\Hess}[1]{\nabla^2 #1}

% -------------------------------------------------------------------------
% Optimization problem objects
% -------------------------------------------------------------------------
\newcommand{\Preg}{P_{\mathrm{reg}}}  % regular perimeter
\newcommand{\Pst}{P_{\mathrm{st}}}   % Steiner perimeter
\newcommand{\Areg}{A^{\mathrm{reg}}} % regular cell area
\newcommand{\Ast}{A^{\mathrm{st}}}   % Steiner area contribution
\newcommand{\Atgt}{\bar{A}}           % target area per cell
\newcommand{\Lagr}{\mathcal{L}}       % Lagrangian
\newcommand{\objfac}{\sigma}          % IPOPT objective scaling factor

% -------------------------------------------------------------------------
% Code cross-references (for "In the code..." callouts)
% -------------------------------------------------------------------------
\newcommand{\code}[1]{\texttt{#1}}
\newcommand{\codefile}[1]{\texttt{\small #1}}

% -------------------------------------------------------------------------
% Threshold dynamics / approach B (10-mbo-auction-dynamics)
% -------------------------------------------------------------------------
\newcommand{\Dlump}{D}                 % lumped mass matrix diag(v)
\newcommand{\Atau}{A_\tau}             % discrete diffusion operator (M + tau K)^{-1} M
\newcommand{\Etau}{E_\tau}             % heat-content (threshold-dynamics) energy
\newcommand{\Rcell}{R_{\mathrm{cell}}} % radius of the equal-area disc sqrt(Abar/pi)
\newcommand{\hmean}{h_{\mathrm{mean}}} % mean edge length
\newcommand{\hmax}{h_{\mathrm{max}}}   % max edge length
\newcommand{\Gt}{G_t}                  % heat kernel at time t
\newcommand{\Ktau}{K^{\mathrm{res}}_\tau} % resolvent (backward-Euler) kernel


% -------------------------------------------------------------------------
% Packages (include here so each document only needs % shared/macros.tex
% Shared notation for all surface-partition math documents.
% Include with: \input{../shared/macros}

% -------------------------------------------------------------------------
% Packages (include here so each document only needs \input{../shared/macros})
% -------------------------------------------------------------------------
\usepackage{amsmath,amssymb,amsthm}
\usepackage{bm}           % \bm for bold math
\usepackage{mathtools}    % \coloneqq, dcases, etc.
\usepackage{booktabs}     % \toprule etc. in tables
\usepackage{hyperref}
\usepackage{geometry}
\usepackage{microtype}
\usepackage{parskip}
\usepackage{xcolor}
\usepackage{tcolorbox}
\usepackage{enumitem}
\usepackage{float}             % [H] float placement

\geometry{a4paper, margin=2.5cm}

% -------------------------------------------------------------------------
% Theorem-like environments
% -------------------------------------------------------------------------
\theoremstyle{definition}
\newtheorem{defn}{Definition}[section]
\newtheorem{remark}{Remark}[section]

\theoremstyle{plain}
\newtheorem{prop}{Proposition}[section]

% -------------------------------------------------------------------------
% Mesh and geometry
% -------------------------------------------------------------------------
\newcommand{\V}{\mathbf{V}}            % vertex array V[i] in R^dim
\newcommand{\F}{\mathbf{F}}            % face (triangle) array
\newcommand{\vone}[1]{v_{1,#1}}       % first vertex index of VP i
\newcommand{\vtwo}[1]{v_{2,#1}}       % second vertex index of VP i

% -------------------------------------------------------------------------
% Variable points
% -------------------------------------------------------------------------
\newcommand{\lam}{\lambda}             % scalar VP parameter
\newcommand{\Lam}{\boldsymbol{\lambda}}% full parameter vector
\newcommand{\vp}[1]{\mathbf{p}_{#1}}  % 3D position of VP i
\newcommand{\ed}[1]{\mathbf{d}_{#1}}  % edge direction d_i = V[v1_i] - V[v2_i]

% -------------------------------------------------------------------------
% Steiner / triple-point
% -------------------------------------------------------------------------
\newcommand{\Sv}{\mathbf{S}}           % Steiner (Fermat-Torricelli) point
\newcommand{\Ki}{K_{i}}               % projection matrix (I-n_i n_i^T)/r_i
\newcommand{\Mmat}{M}                 % sum of K_i matrices at a triple point

% -------------------------------------------------------------------------
% Normals and unit vectors
% -------------------------------------------------------------------------
\newcommand{\nhat}{\hat{\mathbf{n}}}  % unit normal to a triangle
\newcommand{\Dhat}{\hat{\boldsymbol{\Delta}}}  % unit segment direction

% -------------------------------------------------------------------------
% Operations
% -------------------------------------------------------------------------
\newcommand{\norm}[1]{\left\|#1\right\|}
\newcommand{\abs}[1]{\left|#1\right|}
\newcommand{\ip}[2]{\langle #1,\, #2 \rangle}

% Projection perpendicular to a unit vector \uhat
\newcommand{\Proj}[1]{\bigl(\mathbf{I} - #1 #1^\top\bigr)}

% argmax operator (winner-take-all assignment in Phase 1 documents)
\DeclareMathOperator*{\argmax}{arg\,max}
\DeclareMathOperator*{\argmin}{arg\,min}

% -------------------------------------------------------------------------
% Calculus shorthands
% -------------------------------------------------------------------------
\newcommand{\pd}[2]{\dfrac{\partial #1}{\partial #2}}
\newcommand{\pdd}[3]{\dfrac{\partial^2 #1}{\partial #2\,\partial #3}}
\newcommand{\grad}{\nabla}
\newcommand{\Hess}[1]{\nabla^2 #1}

% -------------------------------------------------------------------------
% Optimization problem objects
% -------------------------------------------------------------------------
\newcommand{\Preg}{P_{\mathrm{reg}}}  % regular perimeter
\newcommand{\Pst}{P_{\mathrm{st}}}   % Steiner perimeter
\newcommand{\Areg}{A^{\mathrm{reg}}} % regular cell area
\newcommand{\Ast}{A^{\mathrm{st}}}   % Steiner area contribution
\newcommand{\Atgt}{\bar{A}}           % target area per cell
\newcommand{\Lagr}{\mathcal{L}}       % Lagrangian
\newcommand{\objfac}{\sigma}          % IPOPT objective scaling factor

% -------------------------------------------------------------------------
% Code cross-references (for "In the code..." callouts)
% -------------------------------------------------------------------------
\newcommand{\code}[1]{\texttt{#1}}
\newcommand{\codefile}[1]{\texttt{\small #1}}

% -------------------------------------------------------------------------
% Threshold dynamics / approach B (10-mbo-auction-dynamics)
% -------------------------------------------------------------------------
\newcommand{\Dlump}{D}                 % lumped mass matrix diag(v)
\newcommand{\Atau}{A_\tau}             % discrete diffusion operator (M + tau K)^{-1} M
\newcommand{\Etau}{E_\tau}             % heat-content (threshold-dynamics) energy
\newcommand{\Rcell}{R_{\mathrm{cell}}} % radius of the equal-area disc sqrt(Abar/pi)
\newcommand{\hmean}{h_{\mathrm{mean}}} % mean edge length
\newcommand{\hmax}{h_{\mathrm{max}}}   % max edge length
\newcommand{\Gt}{G_t}                  % heat kernel at time t
\newcommand{\Ktau}{K^{\mathrm{res}}_\tau} % resolvent (backward-Euler) kernel
)
% -------------------------------------------------------------------------
\usepackage{amsmath,amssymb,amsthm}
\usepackage{bm}           % \bm for bold math
\usepackage{mathtools}    % \coloneqq, dcases, etc.
\usepackage{booktabs}     % \toprule etc. in tables
\usepackage{hyperref}
\usepackage{geometry}
\usepackage{microtype}
\usepackage{parskip}
\usepackage{xcolor}
\usepackage{tcolorbox}
\usepackage{enumitem}
\usepackage{float}             % [H] float placement

\geometry{a4paper, margin=2.5cm}

% -------------------------------------------------------------------------
% Theorem-like environments
% -------------------------------------------------------------------------
\theoremstyle{definition}
\newtheorem{defn}{Definition}[section]
\newtheorem{remark}{Remark}[section]

\theoremstyle{plain}
\newtheorem{prop}{Proposition}[section]

% -------------------------------------------------------------------------
% Mesh and geometry
% -------------------------------------------------------------------------
\newcommand{\V}{\mathbf{V}}            % vertex array V[i] in R^dim
\newcommand{\F}{\mathbf{F}}            % face (triangle) array
\newcommand{\vone}[1]{v_{1,#1}}       % first vertex index of VP i
\newcommand{\vtwo}[1]{v_{2,#1}}       % second vertex index of VP i

% -------------------------------------------------------------------------
% Variable points
% -------------------------------------------------------------------------
\newcommand{\lam}{\lambda}             % scalar VP parameter
\newcommand{\Lam}{\boldsymbol{\lambda}}% full parameter vector
\newcommand{\vp}[1]{\mathbf{p}_{#1}}  % 3D position of VP i
\newcommand{\ed}[1]{\mathbf{d}_{#1}}  % edge direction d_i = V[v1_i] - V[v2_i]

% -------------------------------------------------------------------------
% Steiner / triple-point
% -------------------------------------------------------------------------
\newcommand{\Sv}{\mathbf{S}}           % Steiner (Fermat-Torricelli) point
\newcommand{\Ki}{K_{i}}               % projection matrix (I-n_i n_i^T)/r_i
\newcommand{\Mmat}{M}                 % sum of K_i matrices at a triple point

% -------------------------------------------------------------------------
% Normals and unit vectors
% -------------------------------------------------------------------------
\newcommand{\nhat}{\hat{\mathbf{n}}}  % unit normal to a triangle
\newcommand{\Dhat}{\hat{\boldsymbol{\Delta}}}  % unit segment direction

% -------------------------------------------------------------------------
% Operations
% -------------------------------------------------------------------------
\newcommand{\norm}[1]{\left\|#1\right\|}
\newcommand{\abs}[1]{\left|#1\right|}
\newcommand{\ip}[2]{\langle #1,\, #2 \rangle}

% Projection perpendicular to a unit vector \uhat
\newcommand{\Proj}[1]{\bigl(\mathbf{I} - #1 #1^\top\bigr)}

% argmax operator (winner-take-all assignment in Phase 1 documents)
\DeclareMathOperator*{\argmax}{arg\,max}
\DeclareMathOperator*{\argmin}{arg\,min}

% -------------------------------------------------------------------------
% Calculus shorthands
% -------------------------------------------------------------------------
\newcommand{\pd}[2]{\dfrac{\partial #1}{\partial #2}}
\newcommand{\pdd}[3]{\dfrac{\partial^2 #1}{\partial #2\,\partial #3}}
\newcommand{\grad}{\nabla}
\newcommand{\Hess}[1]{\nabla^2 #1}

% -------------------------------------------------------------------------
% Optimization problem objects
% -------------------------------------------------------------------------
\newcommand{\Preg}{P_{\mathrm{reg}}}  % regular perimeter
\newcommand{\Pst}{P_{\mathrm{st}}}   % Steiner perimeter
\newcommand{\Areg}{A^{\mathrm{reg}}} % regular cell area
\newcommand{\Ast}{A^{\mathrm{st}}}   % Steiner area contribution
\newcommand{\Atgt}{\bar{A}}           % target area per cell
\newcommand{\Lagr}{\mathcal{L}}       % Lagrangian
\newcommand{\objfac}{\sigma}          % IPOPT objective scaling factor

% -------------------------------------------------------------------------
% Code cross-references (for "In the code..." callouts)
% -------------------------------------------------------------------------
\newcommand{\code}[1]{\texttt{#1}}
\newcommand{\codefile}[1]{\texttt{\small #1}}

% -------------------------------------------------------------------------
% Threshold dynamics / approach B (10-mbo-auction-dynamics)
% -------------------------------------------------------------------------
\newcommand{\Dlump}{D}                 % lumped mass matrix diag(v)
\newcommand{\Atau}{A_\tau}             % discrete diffusion operator (M + tau K)^{-1} M
\newcommand{\Etau}{E_\tau}             % heat-content (threshold-dynamics) energy
\newcommand{\Rcell}{R_{\mathrm{cell}}} % radius of the equal-area disc sqrt(Abar/pi)
\newcommand{\hmean}{h_{\mathrm{mean}}} % mean edge length
\newcommand{\hmax}{h_{\mathrm{max}}}   % max edge length
\newcommand{\Gt}{G_t}                  % heat kernel at time t
\newcommand{\Ktau}{K^{\mathrm{res}}_\tau} % resolvent (backward-Euler) kernel


% -------------------------------------------------------------------------
% Packages (include here so each document only needs % shared/macros.tex
% Shared notation for all surface-partition math documents.
% Include with: % shared/macros.tex
% Shared notation for all surface-partition math documents.
% Include with: \input{../shared/macros}

% -------------------------------------------------------------------------
% Packages (include here so each document only needs \input{../shared/macros})
% -------------------------------------------------------------------------
\usepackage{amsmath,amssymb,amsthm}
\usepackage{bm}           % \bm for bold math
\usepackage{mathtools}    % \coloneqq, dcases, etc.
\usepackage{booktabs}     % \toprule etc. in tables
\usepackage{hyperref}
\usepackage{geometry}
\usepackage{microtype}
\usepackage{parskip}
\usepackage{xcolor}
\usepackage{tcolorbox}
\usepackage{enumitem}
\usepackage{float}             % [H] float placement

\geometry{a4paper, margin=2.5cm}

% -------------------------------------------------------------------------
% Theorem-like environments
% -------------------------------------------------------------------------
\theoremstyle{definition}
\newtheorem{defn}{Definition}[section]
\newtheorem{remark}{Remark}[section]

\theoremstyle{plain}
\newtheorem{prop}{Proposition}[section]

% -------------------------------------------------------------------------
% Mesh and geometry
% -------------------------------------------------------------------------
\newcommand{\V}{\mathbf{V}}            % vertex array V[i] in R^dim
\newcommand{\F}{\mathbf{F}}            % face (triangle) array
\newcommand{\vone}[1]{v_{1,#1}}       % first vertex index of VP i
\newcommand{\vtwo}[1]{v_{2,#1}}       % second vertex index of VP i

% -------------------------------------------------------------------------
% Variable points
% -------------------------------------------------------------------------
\newcommand{\lam}{\lambda}             % scalar VP parameter
\newcommand{\Lam}{\boldsymbol{\lambda}}% full parameter vector
\newcommand{\vp}[1]{\mathbf{p}_{#1}}  % 3D position of VP i
\newcommand{\ed}[1]{\mathbf{d}_{#1}}  % edge direction d_i = V[v1_i] - V[v2_i]

% -------------------------------------------------------------------------
% Steiner / triple-point
% -------------------------------------------------------------------------
\newcommand{\Sv}{\mathbf{S}}           % Steiner (Fermat-Torricelli) point
\newcommand{\Ki}{K_{i}}               % projection matrix (I-n_i n_i^T)/r_i
\newcommand{\Mmat}{M}                 % sum of K_i matrices at a triple point

% -------------------------------------------------------------------------
% Normals and unit vectors
% -------------------------------------------------------------------------
\newcommand{\nhat}{\hat{\mathbf{n}}}  % unit normal to a triangle
\newcommand{\Dhat}{\hat{\boldsymbol{\Delta}}}  % unit segment direction

% -------------------------------------------------------------------------
% Operations
% -------------------------------------------------------------------------
\newcommand{\norm}[1]{\left\|#1\right\|}
\newcommand{\abs}[1]{\left|#1\right|}
\newcommand{\ip}[2]{\langle #1,\, #2 \rangle}

% Projection perpendicular to a unit vector \uhat
\newcommand{\Proj}[1]{\bigl(\mathbf{I} - #1 #1^\top\bigr)}

% argmax operator (winner-take-all assignment in Phase 1 documents)
\DeclareMathOperator*{\argmax}{arg\,max}
\DeclareMathOperator*{\argmin}{arg\,min}

% -------------------------------------------------------------------------
% Calculus shorthands
% -------------------------------------------------------------------------
\newcommand{\pd}[2]{\dfrac{\partial #1}{\partial #2}}
\newcommand{\pdd}[3]{\dfrac{\partial^2 #1}{\partial #2\,\partial #3}}
\newcommand{\grad}{\nabla}
\newcommand{\Hess}[1]{\nabla^2 #1}

% -------------------------------------------------------------------------
% Optimization problem objects
% -------------------------------------------------------------------------
\newcommand{\Preg}{P_{\mathrm{reg}}}  % regular perimeter
\newcommand{\Pst}{P_{\mathrm{st}}}   % Steiner perimeter
\newcommand{\Areg}{A^{\mathrm{reg}}} % regular cell area
\newcommand{\Ast}{A^{\mathrm{st}}}   % Steiner area contribution
\newcommand{\Atgt}{\bar{A}}           % target area per cell
\newcommand{\Lagr}{\mathcal{L}}       % Lagrangian
\newcommand{\objfac}{\sigma}          % IPOPT objective scaling factor

% -------------------------------------------------------------------------
% Code cross-references (for "In the code..." callouts)
% -------------------------------------------------------------------------
\newcommand{\code}[1]{\texttt{#1}}
\newcommand{\codefile}[1]{\texttt{\small #1}}

% -------------------------------------------------------------------------
% Threshold dynamics / approach B (10-mbo-auction-dynamics)
% -------------------------------------------------------------------------
\newcommand{\Dlump}{D}                 % lumped mass matrix diag(v)
\newcommand{\Atau}{A_\tau}             % discrete diffusion operator (M + tau K)^{-1} M
\newcommand{\Etau}{E_\tau}             % heat-content (threshold-dynamics) energy
\newcommand{\Rcell}{R_{\mathrm{cell}}} % radius of the equal-area disc sqrt(Abar/pi)
\newcommand{\hmean}{h_{\mathrm{mean}}} % mean edge length
\newcommand{\hmax}{h_{\mathrm{max}}}   % max edge length
\newcommand{\Gt}{G_t}                  % heat kernel at time t
\newcommand{\Ktau}{K^{\mathrm{res}}_\tau} % resolvent (backward-Euler) kernel


% -------------------------------------------------------------------------
% Packages (include here so each document only needs % shared/macros.tex
% Shared notation for all surface-partition math documents.
% Include with: \input{../shared/macros}

% -------------------------------------------------------------------------
% Packages (include here so each document only needs \input{../shared/macros})
% -------------------------------------------------------------------------
\usepackage{amsmath,amssymb,amsthm}
\usepackage{bm}           % \bm for bold math
\usepackage{mathtools}    % \coloneqq, dcases, etc.
\usepackage{booktabs}     % \toprule etc. in tables
\usepackage{hyperref}
\usepackage{geometry}
\usepackage{microtype}
\usepackage{parskip}
\usepackage{xcolor}
\usepackage{tcolorbox}
\usepackage{enumitem}
\usepackage{float}             % [H] float placement

\geometry{a4paper, margin=2.5cm}

% -------------------------------------------------------------------------
% Theorem-like environments
% -------------------------------------------------------------------------
\theoremstyle{definition}
\newtheorem{defn}{Definition}[section]
\newtheorem{remark}{Remark}[section]

\theoremstyle{plain}
\newtheorem{prop}{Proposition}[section]

% -------------------------------------------------------------------------
% Mesh and geometry
% -------------------------------------------------------------------------
\newcommand{\V}{\mathbf{V}}            % vertex array V[i] in R^dim
\newcommand{\F}{\mathbf{F}}            % face (triangle) array
\newcommand{\vone}[1]{v_{1,#1}}       % first vertex index of VP i
\newcommand{\vtwo}[1]{v_{2,#1}}       % second vertex index of VP i

% -------------------------------------------------------------------------
% Variable points
% -------------------------------------------------------------------------
\newcommand{\lam}{\lambda}             % scalar VP parameter
\newcommand{\Lam}{\boldsymbol{\lambda}}% full parameter vector
\newcommand{\vp}[1]{\mathbf{p}_{#1}}  % 3D position of VP i
\newcommand{\ed}[1]{\mathbf{d}_{#1}}  % edge direction d_i = V[v1_i] - V[v2_i]

% -------------------------------------------------------------------------
% Steiner / triple-point
% -------------------------------------------------------------------------
\newcommand{\Sv}{\mathbf{S}}           % Steiner (Fermat-Torricelli) point
\newcommand{\Ki}{K_{i}}               % projection matrix (I-n_i n_i^T)/r_i
\newcommand{\Mmat}{M}                 % sum of K_i matrices at a triple point

% -------------------------------------------------------------------------
% Normals and unit vectors
% -------------------------------------------------------------------------
\newcommand{\nhat}{\hat{\mathbf{n}}}  % unit normal to a triangle
\newcommand{\Dhat}{\hat{\boldsymbol{\Delta}}}  % unit segment direction

% -------------------------------------------------------------------------
% Operations
% -------------------------------------------------------------------------
\newcommand{\norm}[1]{\left\|#1\right\|}
\newcommand{\abs}[1]{\left|#1\right|}
\newcommand{\ip}[2]{\langle #1,\, #2 \rangle}

% Projection perpendicular to a unit vector \uhat
\newcommand{\Proj}[1]{\bigl(\mathbf{I} - #1 #1^\top\bigr)}

% argmax operator (winner-take-all assignment in Phase 1 documents)
\DeclareMathOperator*{\argmax}{arg\,max}
\DeclareMathOperator*{\argmin}{arg\,min}

% -------------------------------------------------------------------------
% Calculus shorthands
% -------------------------------------------------------------------------
\newcommand{\pd}[2]{\dfrac{\partial #1}{\partial #2}}
\newcommand{\pdd}[3]{\dfrac{\partial^2 #1}{\partial #2\,\partial #3}}
\newcommand{\grad}{\nabla}
\newcommand{\Hess}[1]{\nabla^2 #1}

% -------------------------------------------------------------------------
% Optimization problem objects
% -------------------------------------------------------------------------
\newcommand{\Preg}{P_{\mathrm{reg}}}  % regular perimeter
\newcommand{\Pst}{P_{\mathrm{st}}}   % Steiner perimeter
\newcommand{\Areg}{A^{\mathrm{reg}}} % regular cell area
\newcommand{\Ast}{A^{\mathrm{st}}}   % Steiner area contribution
\newcommand{\Atgt}{\bar{A}}           % target area per cell
\newcommand{\Lagr}{\mathcal{L}}       % Lagrangian
\newcommand{\objfac}{\sigma}          % IPOPT objective scaling factor

% -------------------------------------------------------------------------
% Code cross-references (for "In the code..." callouts)
% -------------------------------------------------------------------------
\newcommand{\code}[1]{\texttt{#1}}
\newcommand{\codefile}[1]{\texttt{\small #1}}

% -------------------------------------------------------------------------
% Threshold dynamics / approach B (10-mbo-auction-dynamics)
% -------------------------------------------------------------------------
\newcommand{\Dlump}{D}                 % lumped mass matrix diag(v)
\newcommand{\Atau}{A_\tau}             % discrete diffusion operator (M + tau K)^{-1} M
\newcommand{\Etau}{E_\tau}             % heat-content (threshold-dynamics) energy
\newcommand{\Rcell}{R_{\mathrm{cell}}} % radius of the equal-area disc sqrt(Abar/pi)
\newcommand{\hmean}{h_{\mathrm{mean}}} % mean edge length
\newcommand{\hmax}{h_{\mathrm{max}}}   % max edge length
\newcommand{\Gt}{G_t}                  % heat kernel at time t
\newcommand{\Ktau}{K^{\mathrm{res}}_\tau} % resolvent (backward-Euler) kernel
)
% -------------------------------------------------------------------------
\usepackage{amsmath,amssymb,amsthm}
\usepackage{bm}           % \bm for bold math
\usepackage{mathtools}    % \coloneqq, dcases, etc.
\usepackage{booktabs}     % \toprule etc. in tables
\usepackage{hyperref}
\usepackage{geometry}
\usepackage{microtype}
\usepackage{parskip}
\usepackage{xcolor}
\usepackage{tcolorbox}
\usepackage{enumitem}
\usepackage{float}             % [H] float placement

\geometry{a4paper, margin=2.5cm}

% -------------------------------------------------------------------------
% Theorem-like environments
% -------------------------------------------------------------------------
\theoremstyle{definition}
\newtheorem{defn}{Definition}[section]
\newtheorem{remark}{Remark}[section]

\theoremstyle{plain}
\newtheorem{prop}{Proposition}[section]

% -------------------------------------------------------------------------
% Mesh and geometry
% -------------------------------------------------------------------------
\newcommand{\V}{\mathbf{V}}            % vertex array V[i] in R^dim
\newcommand{\F}{\mathbf{F}}            % face (triangle) array
\newcommand{\vone}[1]{v_{1,#1}}       % first vertex index of VP i
\newcommand{\vtwo}[1]{v_{2,#1}}       % second vertex index of VP i

% -------------------------------------------------------------------------
% Variable points
% -------------------------------------------------------------------------
\newcommand{\lam}{\lambda}             % scalar VP parameter
\newcommand{\Lam}{\boldsymbol{\lambda}}% full parameter vector
\newcommand{\vp}[1]{\mathbf{p}_{#1}}  % 3D position of VP i
\newcommand{\ed}[1]{\mathbf{d}_{#1}}  % edge direction d_i = V[v1_i] - V[v2_i]

% -------------------------------------------------------------------------
% Steiner / triple-point
% -------------------------------------------------------------------------
\newcommand{\Sv}{\mathbf{S}}           % Steiner (Fermat-Torricelli) point
\newcommand{\Ki}{K_{i}}               % projection matrix (I-n_i n_i^T)/r_i
\newcommand{\Mmat}{M}                 % sum of K_i matrices at a triple point

% -------------------------------------------------------------------------
% Normals and unit vectors
% -------------------------------------------------------------------------
\newcommand{\nhat}{\hat{\mathbf{n}}}  % unit normal to a triangle
\newcommand{\Dhat}{\hat{\boldsymbol{\Delta}}}  % unit segment direction

% -------------------------------------------------------------------------
% Operations
% -------------------------------------------------------------------------
\newcommand{\norm}[1]{\left\|#1\right\|}
\newcommand{\abs}[1]{\left|#1\right|}
\newcommand{\ip}[2]{\langle #1,\, #2 \rangle}

% Projection perpendicular to a unit vector \uhat
\newcommand{\Proj}[1]{\bigl(\mathbf{I} - #1 #1^\top\bigr)}

% argmax operator (winner-take-all assignment in Phase 1 documents)
\DeclareMathOperator*{\argmax}{arg\,max}
\DeclareMathOperator*{\argmin}{arg\,min}

% -------------------------------------------------------------------------
% Calculus shorthands
% -------------------------------------------------------------------------
\newcommand{\pd}[2]{\dfrac{\partial #1}{\partial #2}}
\newcommand{\pdd}[3]{\dfrac{\partial^2 #1}{\partial #2\,\partial #3}}
\newcommand{\grad}{\nabla}
\newcommand{\Hess}[1]{\nabla^2 #1}

% -------------------------------------------------------------------------
% Optimization problem objects
% -------------------------------------------------------------------------
\newcommand{\Preg}{P_{\mathrm{reg}}}  % regular perimeter
\newcommand{\Pst}{P_{\mathrm{st}}}   % Steiner perimeter
\newcommand{\Areg}{A^{\mathrm{reg}}} % regular cell area
\newcommand{\Ast}{A^{\mathrm{st}}}   % Steiner area contribution
\newcommand{\Atgt}{\bar{A}}           % target area per cell
\newcommand{\Lagr}{\mathcal{L}}       % Lagrangian
\newcommand{\objfac}{\sigma}          % IPOPT objective scaling factor

% -------------------------------------------------------------------------
% Code cross-references (for "In the code..." callouts)
% -------------------------------------------------------------------------
\newcommand{\code}[1]{\texttt{#1}}
\newcommand{\codefile}[1]{\texttt{\small #1}}

% -------------------------------------------------------------------------
% Threshold dynamics / approach B (10-mbo-auction-dynamics)
% -------------------------------------------------------------------------
\newcommand{\Dlump}{D}                 % lumped mass matrix diag(v)
\newcommand{\Atau}{A_\tau}             % discrete diffusion operator (M + tau K)^{-1} M
\newcommand{\Etau}{E_\tau}             % heat-content (threshold-dynamics) energy
\newcommand{\Rcell}{R_{\mathrm{cell}}} % radius of the equal-area disc sqrt(Abar/pi)
\newcommand{\hmean}{h_{\mathrm{mean}}} % mean edge length
\newcommand{\hmax}{h_{\mathrm{max}}}   % max edge length
\newcommand{\Gt}{G_t}                  % heat kernel at time t
\newcommand{\Ktau}{K^{\mathrm{res}}_\tau} % resolvent (backward-Euler) kernel
)
% -------------------------------------------------------------------------
\usepackage{amsmath,amssymb,amsthm}
\usepackage{bm}           % \bm for bold math
\usepackage{mathtools}    % \coloneqq, dcases, etc.
\usepackage{booktabs}     % \toprule etc. in tables
\usepackage{hyperref}
\usepackage{geometry}
\usepackage{microtype}
\usepackage{parskip}
\usepackage{xcolor}
\usepackage{tcolorbox}
\usepackage{enumitem}
\usepackage{float}             % [H] float placement

\geometry{a4paper, margin=2.5cm}

% -------------------------------------------------------------------------
% Theorem-like environments
% -------------------------------------------------------------------------
\theoremstyle{definition}
\newtheorem{defn}{Definition}[section]
\newtheorem{remark}{Remark}[section]

\theoremstyle{plain}
\newtheorem{prop}{Proposition}[section]

% -------------------------------------------------------------------------
% Mesh and geometry
% -------------------------------------------------------------------------
\newcommand{\V}{\mathbf{V}}            % vertex array V[i] in R^dim
\newcommand{\F}{\mathbf{F}}            % face (triangle) array
\newcommand{\vone}[1]{v_{1,#1}}       % first vertex index of VP i
\newcommand{\vtwo}[1]{v_{2,#1}}       % second vertex index of VP i

% -------------------------------------------------------------------------
% Variable points
% -------------------------------------------------------------------------
\newcommand{\lam}{\lambda}             % scalar VP parameter
\newcommand{\Lam}{\boldsymbol{\lambda}}% full parameter vector
\newcommand{\vp}[1]{\mathbf{p}_{#1}}  % 3D position of VP i
\newcommand{\ed}[1]{\mathbf{d}_{#1}}  % edge direction d_i = V[v1_i] - V[v2_i]

% -------------------------------------------------------------------------
% Steiner / triple-point
% -------------------------------------------------------------------------
\newcommand{\Sv}{\mathbf{S}}           % Steiner (Fermat-Torricelli) point
\newcommand{\Ki}{K_{i}}               % projection matrix (I-n_i n_i^T)/r_i
\newcommand{\Mmat}{M}                 % sum of K_i matrices at a triple point

% -------------------------------------------------------------------------
% Normals and unit vectors
% -------------------------------------------------------------------------
\newcommand{\nhat}{\hat{\mathbf{n}}}  % unit normal to a triangle
\newcommand{\Dhat}{\hat{\boldsymbol{\Delta}}}  % unit segment direction

% -------------------------------------------------------------------------
% Operations
% -------------------------------------------------------------------------
\newcommand{\norm}[1]{\left\|#1\right\|}
\newcommand{\abs}[1]{\left|#1\right|}
\newcommand{\ip}[2]{\langle #1,\, #2 \rangle}

% Projection perpendicular to a unit vector \uhat
\newcommand{\Proj}[1]{\bigl(\mathbf{I} - #1 #1^\top\bigr)}

% argmax operator (winner-take-all assignment in Phase 1 documents)
\DeclareMathOperator*{\argmax}{arg\,max}
\DeclareMathOperator*{\argmin}{arg\,min}

% -------------------------------------------------------------------------
% Calculus shorthands
% -------------------------------------------------------------------------
\newcommand{\pd}[2]{\dfrac{\partial #1}{\partial #2}}
\newcommand{\pdd}[3]{\dfrac{\partial^2 #1}{\partial #2\,\partial #3}}
\newcommand{\grad}{\nabla}
\newcommand{\Hess}[1]{\nabla^2 #1}

% -------------------------------------------------------------------------
% Optimization problem objects
% -------------------------------------------------------------------------
\newcommand{\Preg}{P_{\mathrm{reg}}}  % regular perimeter
\newcommand{\Pst}{P_{\mathrm{st}}}   % Steiner perimeter
\newcommand{\Areg}{A^{\mathrm{reg}}} % regular cell area
\newcommand{\Ast}{A^{\mathrm{st}}}   % Steiner area contribution
\newcommand{\Atgt}{\bar{A}}           % target area per cell
\newcommand{\Lagr}{\mathcal{L}}       % Lagrangian
\newcommand{\objfac}{\sigma}          % IPOPT objective scaling factor

% -------------------------------------------------------------------------
% Code cross-references (for "In the code..." callouts)
% -------------------------------------------------------------------------
\newcommand{\code}[1]{\texttt{#1}}
\newcommand{\codefile}[1]{\texttt{\small #1}}

% -------------------------------------------------------------------------
% Threshold dynamics / approach B (10-mbo-auction-dynamics)
% -------------------------------------------------------------------------
\newcommand{\Dlump}{D}                 % lumped mass matrix diag(v)
\newcommand{\Atau}{A_\tau}             % discrete diffusion operator (M + tau K)^{-1} M
\newcommand{\Etau}{E_\tau}             % heat-content (threshold-dynamics) energy
\newcommand{\Rcell}{R_{\mathrm{cell}}} % radius of the equal-area disc sqrt(Abar/pi)
\newcommand{\hmean}{h_{\mathrm{mean}}} % mean edge length
\newcommand{\hmax}{h_{\mathrm{max}}}   % max edge length
\newcommand{\Gt}{G_t}                  % heat kernel at time t
\newcommand{\Ktau}{K^{\mathrm{res}}_\tau} % resolvent (backward-Euler) kernel
      % loads ALL packages and notation

\hypersetup{
  colorlinks = true,
  linkcolor  = blue!60!black,
  citecolor  = green!50!black,
  urlcolor   = blue!60!black,
  pdftitle   = {The Phase 1 constraint-projection QP and its concave dual},
  pdfauthor  = {surface-partition project}
}

\title{\textbf{The Phase~1 Constraint-Projection QP and its Concave Dual}\\[0.5em]
  \large Exact Euclidean projection onto partition, equal-area, and box
  constraints: dual derivation, the per-vertex simplex solve, the outer
  Jacobian, and the exactness/idempotency guarantees}
\author{}
\date{July 2026}

\begin{document}
\maketitle

% ---------------------------------------------------------------------
% STATUS BANNER
% ---------------------------------------------------------------------
\begin{tcolorbox}[colback=orange!8!white, colframe=orange!65!black,
                  title={Status --- specification (pre-implementation)}, fontupper=\small]
This document is the mathematical \emph{source of truth} for the planned function
\code{orthogonal\_projection\_newton} in \codefile{src/optimization/projection.py}
(design: \codefile{docs/plans/PHASE1\_DUAL\_NEWTON\_PROJECTION\_PLAN.md}). Unlike the
other \codefile{docs/math/} documents, it is written \emph{before} the code (the plan's
workflow: derive first, then implement). The code callout boxes therefore name a
function that does not yet exist; they define what it must compute. The measured
motivation --- that the current \code{orthogonal\_projection\_iterative} does \emph{not}
compute this projection --- is documented in
\codefile{docs/experiments/03-dual-projection-verification/}.
\end{tcolorbox}

\begin{abstract}
Each Phase~1 PGD iteration projects a candidate density onto the intersection of the
partition constraint ($\sum_k u_{i,k}=1$ per vertex), the equal-area constraint
($\sum_i v_i u_{i,k}=d_k$ per cell), and the box $0\le u\le1$. This document derives the
exact Euclidean projection through its Lagrangian dual. The primal is recovered from the
duals by a clip (\eqref{eq:primal-from-dual}); eliminating the per-vertex row duals is a
probability-simplex projection, which must be computed \emph{cap-free} for the row dual to
be well defined (\S\ref{sec:inner}). What remains is the maximization of a concave,
$C^1$, piecewise-quadratic \emph{partial dual} $q(\beta)$ of dimension $N$ (the cell
count), with $\grad q=-R$ the negated area residual (\S\ref{sec:qbeta}). We derive its
generalized Hessian $-J$ (the $N\times N$ area Jacobian, \eqref{eq:jacobian}), prove it is
symmetric positive semidefinite with the \emph{structural} kernel $\operatorname{span}\{\mathbf 1\}$
(\S\ref{sec:jacobian}), and use this to justify the solver: globally-convergent L-BFGS on
$q$ (primary), with a gauge-fixed semismooth-Newton polish (\S\ref{sec:solve}). Finally we
prove exactness and idempotency (\S\ref{sec:exact}). All quantities are analytical.
\end{abstract}

\tableofcontents
\newpage

% =====================================================================
\section{Overview}
\label{sec:overview}
% =====================================================================
Phase~1 minimizes a $\Gamma$-convergence energy by projected gradient descent on a nodal
density $U\in\mathbb{R}^{V\times N}$ ($V$ mesh vertices, $N$ cells). Each step takes a
post-gradient candidate $Y$ (generally infeasible) and returns the nearest feasible density
--- the Euclidean projection onto the constraint set. This projection dominates Phase-1 wall
time and, as measured in \codefile{docs/experiments/03-dual-projection-verification/}, the
incumbent iterative routine does not actually compute it. Here we derive the exact
projection and the structure that makes it cheap: a per-vertex closed form plus a small
$N$-dimensional concave maximization.

% =====================================================================
\section{Notation and Setup}
\label{sec:setup}
% =====================================================================
$U=(u_{i,k})\in\mathbb{R}^{V\times N}$ is the density: row $i$ is vertex $i$, column $k$ is
cell $k$. $v\in\mathbb{R}^V$ is the lumped $P_1$ mass (positive), $\mathbf 1$ a vector of
ones (length inferred), $\Delta=\{p\in\mathbb{R}^N: \sum_k p_k=1,\ p\ge0\}$ the probability
simplex. Given the candidate $Y$ and per-cell area targets $d\in\mathbb{R}^N$
($d_k=\tfrac1N\sum_i v_i$ in the equal-area case), the projection is

\begin{equation}
  \min_{U}\ \tfrac12\norm{U-Y}^2
  \quad\text{s.t.}\quad
  \underbrace{\textstyle\sum_k u_{i,k}=1\ \ \forall i}_{\text{row / partition}},\quad
  \underbrace{\textstyle\sum_i v_i u_{i,k}=d_k\ \ \forall k}_{\text{area}},\quad
  0\le u_{i,k}\le 1 .
  \label{eq:qp}
\end{equation}

The metric is the plain (unweighted) Euclidean norm --- verified against the incumbent's
correction term in \codefile{docs/experiments/03-…} and \S\ref{sec:mechanism-note}.

\begin{remark}[the upper box bound is redundant]
\label{rem:cap}
If $\sum_k u_{i,k}=1$ and $u_{i,k}\ge0$ then each $u_{i,k}\le1$ automatically. Hence the
box in \eqref{eq:qp} reduces to $U\ge0$, and the per-vertex feasible set is exactly $\Delta$.
The cap at $1$ never binds; it is nonetheless important \emph{not} to carry it inside the
inner solve (\S\ref{sec:inner}, Remark~\ref{rem:capfree}).
\end{remark}

\begin{remark}[feasibility of \eqref{eq:qp}]
\label{rem:feasible}
Summing the row constraints weighted by $v$ gives $\sum_k\sum_i v_i u_{i,k}=\sum_i v_i$, so
the area targets must satisfy the consistency condition $\sum_k d_k=\sum_i v_i$. Under it
(and $d\ge0$) the constant-column density $u_{i,k}=d_k/\sum_i v_i\in(0,1)$ is strictly
feasible, so \eqref{eq:qp} is feasible and Slater's condition holds --- strong duality
applies.
\end{remark}

% =====================================================================
\section{The dual and the primal-from-dual map}
\label{sec:dual}
% =====================================================================
Introduce a multiplier $\alpha_i$ for each row constraint and $\beta_k$ for each area
constraint. Keeping the box $0\le U\le1$ explicit, the Lagrangian is
\begin{equation}
  \Lagr(U,\alpha,\beta)=\tfrac12\norm{U-Y}^2
   -\sum_i\alpha_i\Bigl(\sum_k u_{i,k}-1\Bigr)
   -\sum_k\beta_k\Bigl(\sum_i v_i u_{i,k}-d_k\Bigr).
  \label{eq:lagrangian}
\end{equation}
The dual function is $g(\alpha,\beta)=\min_{0\le U\le1}\Lagr(U,\alpha,\beta)$. Because
\eqref{eq:lagrangian} separates over entries, the minimizing $U$ is obtained entrywise: the
unconstrained stationary point is $y_{i,k}+\alpha_i+v_i\beta_k$, clipped to the box, giving
the \textbf{primal-from-dual map}
\begin{equation}
  u_{i,k}(\alpha,\beta)=\operatorname{clip}\!\bigl(y_{i,k}+\alpha_i+v_i\beta_k,\ 0,\ 1\bigr).
  \label{eq:primal-from-dual}
\end{equation}
Note the row dual $\alpha_i$ enters with unit weight and the area dual $\beta_k$ with the
mass weight $v_i$ --- the signature of the \emph{Euclidean} (not $v$-weighted) metric.
\label{sec:mechanism-note}
Substituting \eqref{eq:primal-from-dual} into \eqref{eq:lagrangian},
\begin{equation}
  g(\alpha,\beta)=\sum_{i,k}\Bigl[\tfrac12\bigl(u_{i,k}-y_{i,k}\bigr)^2
   -(\alpha_i+v_i\beta_k)\,u_{i,k}\Bigr]
   +\sum_i\alpha_i+\sum_k\beta_k d_k ,
  \label{eq:dualfun}
\end{equation}
a concave function of $(\alpha,\beta)$ (a pointwise minimum, over $U$, of maps that are
affine in $(\alpha,\beta)$). By Danskin's theorem \cite{bertsekas1999nonlinear}, with
$u=u(\alpha,\beta)$ optimal,
\begin{equation}
  \pd{g}{\alpha_i}=1-\sum_k u_{i,k},\qquad
  \pd{g}{\beta_k}=d_k-\sum_i v_i u_{i,k}=:-R_k(\alpha,\beta),
  \label{eq:dualgrad}
\end{equation}
where $R\in\mathbb{R}^N$ is the \emph{area residual}. The projection is $U(\alpha^\star,\beta^\star)$
at the dual maximizer, where both gradients vanish (feasibility) and \eqref{eq:primal-from-dual}
holds (KKT); by Remark~\ref{rem:feasible} strong duality makes this exact.

\begin{tcolorbox}[colback=blue!4!white, colframe=blue!40!black,
                  title={Code: primal reconstruction}, fontupper=\small]
\codefile{src/optimization/projection.py}: \code{orthogonal\_projection\_newton} (planned).\\[2pt]
The final primal is \code{U = clip(Y + alpha[:,None] + v[:,None]*beta[None,:], 0, 1)}
(\eqref{eq:primal-from-dual}); \code{alpha}, \code{beta} come from \S\ref{sec:inner}--\S\ref{sec:solve}.
\end{tcolorbox}

% =====================================================================
\section{Eliminating the row duals: the per-vertex simplex solve}
\label{sec:inner}
% =====================================================================
For fixed $\beta$, \eqref{eq:dualfun} separates over vertices, so each $\alpha_i$ is found
independently by maximizing $g$ in $\alpha_i$, i.e.\ by solving $\partial g/\partial\alpha_i=0$
from \eqref{eq:dualgrad}. Writing $z_{i,k}:=y_{i,k}+v_i\beta_k$,
\begin{equation}
  \sum_k \operatorname{clip}\!\bigl(z_{i,k}+\alpha_i,\ 0,\ 1\bigr)=1 .
  \label{eq:innereq}
\end{equation}
By Remark~\ref{rem:cap} this is the projection of $z_i$ onto the probability simplex
$\Delta$: the minimizer of $\tfrac12\norm{u_i-z_i}^2$ over $u_i\in\Delta$ has the form
$u_{i,k}=\max(z_{i,k}-\tau_i,0)$ with the threshold $\tau_i=-\alpha_i$ fixed by
$\sum_k u_{i,k}=1$. This is Condat's exact finite (sort/threshold) solve \cite{condat2016simplex}.

\begin{prop}[well-posedness of the row dual]
\label{prop:innerunique}
Solve \eqref{eq:innereq} \emph{cap-free}, i.e.\ as projection onto $\Delta=\{\sum=1,\ \cdot\ge0\}$.
Then $\alpha_i$ is the unique root of $h_i(\alpha):=\sum_k\max(z_{i,k}+\alpha,0)=1$.
\end{prop}
\begin{proof}
$h_i$ is continuous, nondecreasing, convex and piecewise linear, with $h_i\to0$ as
$\alpha\to-\infty$ and $h_i\to+\infty$ as $\alpha\to+\infty$. On any interval where at least
one argument is positive its slope equals the number of active entries $\ge1$, so $h_i$ is
\emph{strictly} increasing there. Since the level $1>0$ is attained only where some entry is
active, the root is unique.
\end{proof}

\begin{remark}[cap-free is a well-posedness requirement, not an optimization]
\label{rem:capfree}
If instead the cap is retained inside \eqref{eq:innereq}, $\tilde h_i(\alpha)=\sum_k
\operatorname{clip}(z_{i,k}+\alpha,0,1)$ has \emph{flat} segments wherever entries saturate at
$1$: whenever a row's largest argument exceeds the runner-up by more than $1$, the level $1$
falls on such a plateau and \eqref{eq:innereq} has an interval of solutions in $\alpha_i$.
E.g.\ $z_i=(2.5,-0.3,-0.8)$ gives $\tilde h_i(\alpha)=1$ for every $\alpha\in[-1.5,0.3]$.
Then $\alpha_i(\beta)$ is set-valued and $d\alpha_i/d\beta$ --- needed for the outer Jacobian
(\S\ref{sec:jacobian}) --- is undefined. The cap-free solve returns the unique
$\alpha_i=1-\max_k z_{i,k}$ (here $-1.5$) and the identical primal $u_i=(1,0,0)$
(Remark~\ref{rem:cap}). The $V$ vertices are independent, so the whole pass is a single
vectorized sort/threshold; it is \emph{stateless}, recomputed from $\beta$ each outer step.
\end{remark}

\begin{tcolorbox}[colback=blue!4!white, colframe=blue!40!black,
                  title={Code: inner solve}, fontupper=\small]
\codefile{src/optimization/projection.py}: \code{orthogonal\_projection\_newton} (planned).\\[2pt]
Vectorized cap-free simplex projection of \code{Z = Y + v[:,None]*beta[None,:]} onto
$\Delta$ (sort each row, find $\tau_i$, \code{U = maximum(Z - tau[:,None], 0)});
\code{alpha = -tau}. Never carry \code{clip(.,0,1)} into this solve (Remark~\ref{rem:capfree}).
\end{tcolorbox}

% =====================================================================
\section{The concave partial dual \texorpdfstring{$q(\beta)$}{q(beta)}}
\label{sec:qbeta}
% =====================================================================
Let $\alpha(\beta)$ denote the row duals of \S\ref{sec:inner} and define the \textbf{partial
dual}
\begin{equation}
  q(\beta):=\max_{\alpha}\ g(\alpha,\beta)=g\bigl(\alpha(\beta),\beta\bigr).
  \label{eq:qdef}
\end{equation}

\begin{prop}[$q$ is concave, $C^1$, with $\grad q=-R$]
\label{prop:q}
$q$ is concave and continuously differentiable, and
\begin{equation}
  \grad q(\beta)=-R(\beta),\qquad R_k(\beta)=\sum_i v_i\,u_{i,k}\bigl(\alpha(\beta),\beta\bigr)-d_k ,
  \label{eq:gradq}
\end{equation}
where $u(\alpha(\beta),\beta)$ is the primal \eqref{eq:primal-from-dual}. Maximizing $q$ over
$\beta$ solves \eqref{eq:qp}.
\end{prop}
\begin{proof}
$g$ is jointly concave (\S\ref{sec:dual}); partial maximization over the convex variable
$\alpha$ preserves concavity, so $q$ is concave. The inner maximizer $\alpha(\beta)$ is
unique (Prop.~\ref{prop:innerunique}) and Lipschitz in $\beta$; by Danskin's theorem
\cite{bertsekas1999nonlinear} $q$ is differentiable with $\grad q(\beta)=\partial g/\partial\beta
\big|_{\alpha=\alpha(\beta)}=d-\sum_i v_i u_{i,\cdot}=-R$ (using $\partial g/\partial\alpha=0$
at $\alpha(\beta)$). The residual $R$ is continuous (the simplex projection is $1$-Lipschitz),
so $q\in C^1$. At $\beta^\star=\arg\max q$, $R(\beta^\star)=0$ and the inner solve enforces
the row constraints, so $u(\alpha(\beta^\star),\beta^\star)$ is primal feasible and satisfies
\eqref{eq:primal-from-dual} --- the KKT conditions of \eqref{eq:qp}; strong duality
(Remark~\ref{rem:feasible}) makes it the projection.
\end{proof}

\emph{Consequence for the solver.} The outer problem is the maximization of a concave $C^1$
function of only $N$ variables, whose value and gradient cost one vectorized inner pass. This
is what makes the projection cheap and robust (\S\ref{sec:solve}).

% =====================================================================
\section{The gauge and the structural singularity}
\label{sec:gauge}
% =====================================================================
\begin{prop}[shift invariance / gauge]
\label{prop:gauge}
For any $t\in\mathbb{R}$, $\ (\alpha,\beta)\mapsto(\alpha-t\,v,\ \beta+t\,\mathbf 1)$ leaves
every $u_{i,k}$ in \eqref{eq:primal-from-dual} unchanged. Consequently $q(\beta+t\mathbf 1)=q(\beta)$,
$\grad q\cdot\mathbf 1\equiv0$, and $\mathbf 1^\top R\equiv0$ after the inner solve.
\end{prop}
\begin{proof}
The argument of the clip is $y_{i,k}+\alpha_i+v_i\beta_k$; under the shift it becomes
$y_{i,k}+(\alpha_i-tv_i)+v_i(\beta_k+t)=y_{i,k}+\alpha_i+v_i\beta_k$, unchanged, so $u$ (hence
$g$) is invariant. As $\beta\mapsto\beta+t\mathbf 1$ is matched by $\alpha(\beta+t\mathbf 1)
=\alpha(\beta)-t v$, $q(\beta+t\mathbf 1)=g(\alpha(\beta),\beta)=q(\beta)$; differentiating in
$t$ gives $\grad q\cdot\mathbf 1=0$, i.e.\ $\mathbf 1^\top R=0$. Directly:
$\mathbf 1^\top R=\sum_i v_i\bigl(\sum_k u_{i,k}\bigr)-\sum_k d_k=\sum_i v_i-\sum_k d_k=0$ by
Remark~\ref{rem:feasible} and $\sum_k u_{i,k}=1$.
\end{proof}

Thus the dual maximizer is a line $\beta^\star+\operatorname{span}\{\mathbf 1\}$; the
\emph{primal} $u$ is unique on it. Any Newton-type linear solve must fix this gauge
(e.g.\ impose $\sum_k\beta_k=0$, or drop one component); L-BFGS simply drifts harmlessly along
the ridge. The residual identity $\mathbf 1^\top R=0$ means the Newton system of
\S\ref{sec:jacobian} is always consistent on $\mathbf 1^\perp$.

% =====================================================================
\section{The outer Jacobian \texorpdfstring{$J=-\Hess{q}$}{J = -Hess q}}
\label{sec:jacobian}
% =====================================================================
Away from the finitely many $\beta$ where an entry crosses a box breakpoint, the active sets
\begin{equation}
  A_i=\{k:\ z_{i,k}+\alpha_i>0\},\qquad m_i:=\abs{A_i}\ge1,
\end{equation}
are locally constant, and (cap-free, Prop.~\ref{prop:innerunique})
$u_{i,k}=(z_{i,k}+\alpha_i)\,[k\in A_i]$ with $\alpha_i=(1-\sum_{k\in A_i}z_{i,k})/m_i$. Since
$\partial z_{i,k}/\partial\beta_l=v_i\delta_{kl}$,
\begin{equation}
  \pd{\alpha_i}{\beta_l}=-\frac{v_i}{m_i}\,[l\in A_i],\qquad
  \pd{u_{i,k}}{\beta_l}=\Bigl(v_i\delta_{kl}-\frac{v_i}{m_i}[l\in A_i]\Bigr)[k\in A_i].
  \label{eq:dudbeta}
\end{equation}
Differentiating $R_k=\sum_i v_i u_{i,k}-d_k$ gives the generalized Jacobian
\begin{equation}
  \boxed{\,J_{kl}=\pd{R_k}{\beta_l}
   =\delta_{kl}\!\!\sum_{i:\,k\in A_i}\!\! v_i^2
   \;-\!\!\sum_{i:\,k,l\in A_i}\!\!\frac{v_i^2}{m_i}\,},\qquad
  J=\sum_i v_i^2\Bigl(\operatorname{diag}(a_i)-\tfrac{1}{m_i}a_i a_i^\top\Bigr),
  \label{eq:jacobian}
\end{equation}
with $a_i\in\{0,1\}^N$ the indicator of $A_i$ ($m_i=a_i^\top a_i$). Since $R=-\grad q$,
$J=-\Hess{q}$.

\begin{prop}[structure of $J$]
\label{prop:J}
$J$ is symmetric positive semidefinite, and $J\mathbf 1=0$ \emph{structurally} (for every
active-set configuration). Its kernel is $\{x:\ x\ \text{constant on each }A_i\}$; if the
active sets jointly connect all $N$ cells this is exactly $\operatorname{span}\{\mathbf 1\}$
and $J$ has rank $N-1$.
\end{prop}
\begin{proof}
Each summand $B_i=\operatorname{diag}(a_i)-\tfrac1{m_i}a_ia_i^\top$ is symmetric, and for any
$x$, $x^\top B_i x=\sum_{k\in A_i}x_k^2-\tfrac1{m_i}\bigl(\sum_{k\in A_i}x_k\bigr)^2\ge0$ by
Cauchy--Schwarz, with equality iff $x$ is constant on $A_i$. Hence $J=\sum_i v_i^2 B_i\succeq0$.
Also $B_i a_i^{(\mathbf 1)}$ where $B_i\mathbf 1=a_i-a_i(a_i^\top\mathbf 1)/m_i=a_i-a_i=0$, so
$J\mathbf 1=0$ regardless of the active sets. The kernel is $\bigcap_i\{x\text{ const on }A_i\}$;
connectedness collapses it to the constants.
\end{proof}

The only linear system anywhere in the method is the \emph{gauge-fixed} reduced
$(N-1)\times(N-1)$ solve of \S\ref{sec:solve}; $N$ is the cell count ($\le\!\sim\!10^3$), so
it is a trivial dense factorization.

\begin{tcolorbox}[colback=blue!4!white, colframe=blue!40!black,
                  title={Code: outer gradient / Jacobian}, fontupper=\small]
\codefile{src/optimization/projection.py}: \code{orthogonal\_projection\_newton} (planned).\\[2pt]
Gradient \code{-R} from one inner pass (\eqref{eq:gradq}); Jacobian \eqref{eq:jacobian} only
if the Newton polish is used, assembled from the active-set indicators $a_i$ and solved on
$\{\sum\beta=0\}$ (never an $\varepsilon$-ridge --- see \S\ref{sec:solve}).
\end{tcolorbox}

% =====================================================================
\section{Solving the outer problem}
\label{sec:solve}
% =====================================================================
\paragraph{Primary --- L-BFGS on the concave dual.} Maximize $q$ (equivalently minimize
$-q$) by L-BFGS-B \cite{byrd1995lbfgsb}; each evaluation is one inner pass
(\S\ref{sec:inner}) returning $q$ and $\grad q=-R$. Because $q$ is concave and $C^1$
(Prop.~\ref{prop:q}), L-BFGS with a Wolfe line search converges globally to the maximizer
(unique modulo the $\mathbf 1$-gauge, Prop.~\ref{prop:gauge}); the primal is unique. No
Jacobian, no linear solve, no active-set bookkeeping.

\paragraph{Optional polish --- semismooth Newton on the gauge-fixed quotient.} $R$ is
semismooth (piecewise smooth in $\beta$); the Newton step solves $J\,\delta\beta=-R$ for
$\delta\beta\in\mathbf 1^\perp$ (consistent by $\mathbf 1^\top R=0$). Since $J\succeq0$,
$\delta\beta=-J^{+}R$ is an ascent direction for $q$ ($\grad q\cdot\delta\beta=R^\top J^{+}R\ge0$),
so use $q$ itself as the line-search merit. On the quotient $\{\sum\beta=0\}$, if the active
sets connect all cells, $J$ is positive definite (Prop.~\ref{prop:J}) and the generalized
Jacobians at $\beta^\star$ are nonsingular (BD-regularity), so Qi--Sun
\cite{qi1993nonsmooth,hintermuller2002semismooth} gives local quadratic convergence there. A
handful of Newton steps drive feasibility from the L-BFGS tolerance to machine level.

\begin{remark}[why not an $\varepsilon$-ridge; the empty-cell degeneracy]
\label{rem:deadcell}
The structural kernel $\operatorname{span}\{\mathbf 1\}$ must be removed by the reduced/gauge-fixed
solve, \emph{not} by a Tikhonov ridge $J+\varepsilon I$. The ridge is not merely inelegant: if
some cell $k$ has \emph{no} active entry at any vertex (an empty cell), then row $k$ of $J$ is
zero while $R_k=-d_k\ne0$, so $J\,\delta\beta=-R$ is inconsistent in coordinate $k$; the ridge
returns $\delta\beta_k=d_k/\varepsilon$ (an explosion, $\sim\!10^{10}$ at $\varepsilon=10^{-4}$),
while a least-squares solve leaves $\delta\beta_k=0$ and never revives the cell. The concave
ascent has no such failure: $\partial q/\partial\beta_k=-R_k=d_k>0$ drives $\beta_k$ up until
cell $k$ re-enters some active set. Hence L-BFGS is the robust primary; the Newton polish must
detect empty active columns and fall back to an ascent step there.
\end{remark}

\begin{tcolorbox}[colback=blue!4!white, colframe=blue!40!black,
                  title={Code: outer solve}, fontupper=\small]
\codefile{src/optimization/projection.py}: \code{orthogonal\_projection\_newton} (planned).\\[2pt]
Primary: \code{scipy.optimize.minimize(-q, beta0, jac=..., method="L-BFGS-B")}, warm-started on
$\beta$ from the previous PGD step. Optional Newton polish: gauge-fixed reduced solve of
\eqref{eq:jacobian} with $q$-ascent line search and the empty-cell safeguard
(Remark~\ref{rem:deadcell}). Thread only $\beta$ ($\alpha$ is recomputed).
\end{tcolorbox}

% =====================================================================
\section{Exactness and idempotency}
\label{sec:exact}
% =====================================================================
\begin{prop}[exactness and idempotency]
\label{prop:exact}
At the dual maximizer $\beta^\star$, $U^\star=U(\alpha(\beta^\star),\beta^\star)$ from
\eqref{eq:primal-from-dual} is the exact solution of \eqref{eq:qp}. Moreover the projection is
idempotent: if $Y$ already satisfies \eqref{eq:qp}, then $U^\star=Y$ to machine precision.
\end{prop}
\begin{proof}
\emph{Exactness.} At $\beta^\star$, $\grad q(\beta^\star)=-R(\beta^\star)=0$ (area feasibility)
and the inner solve enforces $\sum_k u^\star_{i,k}=1$ (row feasibility) with $u^\star\ge0$;
\eqref{eq:primal-from-dual} is the box-KKT stationarity with multipliers
$(\alpha(\beta^\star),\beta^\star)$. These are the KKT conditions of the convex program
\eqref{eq:qp}, which are sufficient; strong duality (Remark~\ref{rem:feasible}) identifies
$U^\star$ with the projection. \emph{Idempotency.} If $Y$ is feasible then $\beta=0$,
$\alpha=0$ give $R=0$ and reproduce $Y$ exactly, so $\beta^\star=0$ is a maximizer and
$U^\star=\operatorname{clip}(Y,0,1)=Y$. A warm-started solve terminates at once, with movement
at the level of the solver tolerance.
\end{proof}

This is precisely what the incumbent lacks: it is not the projection at all (gap up to $0.86$;
\codefile{docs/experiments/03-dual-projection-verification/}), and its
$\operatorname{clip}\!\circ\!\text{project}$ step is only weakly idempotent
($\sim\!6\times10^{-8}$ per iteration). The dual projection removes both.

% =====================================================================
\section*{Summary table}
% =====================================================================
\begin{center}\small
\begin{tabular}{@{}lll@{}}
\toprule
Quantity & Form & Where \\
\midrule
Primal from duals & $u_{i,k}=\operatorname{clip}(y_{i,k}+\alpha_i+v_i\beta_k,0,1)$ & \eqref{eq:primal-from-dual}, analytical \\
Row dual $\alpha_i(\beta)$ & cap-free simplex projection of $z_i$ (unique) & \S\ref{sec:inner}, analytical \\
Partial dual $q(\beta)$ & concave, $C^1$, $\grad q=-R$ & Prop.~\ref{prop:q}, analytical \\
Gauge & $(\alpha,\beta)\!\to\!(\alpha-tv,\beta+t\mathbf 1)$; $\grad q\cdot\mathbf 1=0$ & Prop.~\ref{prop:gauge} \\
Outer Jacobian $J=-\Hess q$ & $\delta_{kl}\sum_{i:k\in A_i}v_i^2-\sum_{i:k,l\in A_i}v_i^2/m_i$ & \eqref{eq:jacobian}, analytical \\
$J$ structure & sym.\ PSD, $J\mathbf 1=0$, rank $N-1$ & Prop.~\ref{prop:J} \\
Primary solver & L-BFGS on $q$ (global) & \S\ref{sec:solve} \\
Polish & semismooth Newton on $\{\sum\beta=0\}$ (quadratic) & \S\ref{sec:solve} \\
Guarantee & exact projection; idempotent & Prop.~\ref{prop:exact} \\
\bottomrule
\end{tabular}
\end{center}

\bibliographystyle{plain}
\bibliography{../shared/references}

\end{document}
